\documentclass[10pt]{article}
\usepackage[margin=20mm]{geometry}
\usepackage{amsmath,amssymb,fontspec,unicode-math,xurl,hyperref}
\IfFontExistsTF{DejaVu Serif}{\setmainfont{DejaVu Serif}}{\setmainfont{Latin Modern Roman}}
\setmathfont{latinmodern-math.otf}
\setlength{\parindent}{0pt}\setlength{\parskip}{6pt}\sloppy
\title{Technical Verification of Audio Flamingo 2}
\author{AI-assisted technical review}
\date{25 October 2026}
\begin{document}
\maketitle
Abstract: The original 48-page paper, its appendices and a pinned Audio Flamingo 2 release were inspected. The supported overall assessment is Minor: the printed loader procedures, one ablation aggregate, GPU-hour units, and released conditioning/cache logic contain specific defects. The central contrastive objective and component-level attention arithmetic have no demonstrated mathematical error in the checked scope. LongAudioBench holdout provenance and judge calibration remain Needs Evidence. No historical benchmark failure or empirical score correction was established. Scientific adjudication is pending.\par
\section*{Paper identity and scope}
Paper 0208 / I25-43577 is Audio Flamingo 2: An Audio-Language Model with Long-Audio Understanding and Expert Reasoning Abilities, by Sreyan Ghosh, Zhifeng Kong, Sonal Kumar, S Sakshi, Jaehyeon Kim, Wei Ping, Rafael Valle, Dinesh Manocha and Bryan Catanzaro, ICML 2025 [paper]. All page references are to the original PDF; its printed and PDF page numbers coincide. The inspected repository is `NVIDIA/audio-flamingo`, ref `audio\_flamingo\_2`, commit `7684e9cb48ede9dcf6f6c0a1fbac091702da4bde`, dated 3 August 2025 [code]. This later release is not a verified historical experiment commit. Its README also describes checkpoint update v20250625.\par
The review covers all paper text, principal tables and algorithms, objective and architecture, curriculum and empirical controls, generation/grading prompts, and relevant released data/model/evaluation paths. All 246 preserved archive members were byte-checked. Semantic inspection concentrates on the main loader, model wrappers, attention helpers, configurations and evaluation logic, not all 146 implementation files. Audio, checkpoints, raw predictions and historical training logs were unavailable. No author code was executed.\par
\section*{Definitions, dependencies and positive checks}
AF2 combines dense AF-CLAP features, three audio transformation layers, gated cross-attention and a frozen Qwen2.5 language model. The 3B designation describes the language model; Table 20 discloses 4.7B total parameters. The three-stage curriculum extends context from 3 to 9 to 30 ten-second windows.\par
For audio \(x_i\), let \(S(i,i)\) sum exponentiated similarities over its positive captions; \(S(j,i)\) uses another audio’s positives and \(S_{\rm neg}(i)\) uses its compositionally altered negatives. The loss on p. 4 is \[\mathcal{L}=-B^{-1}\sum_{i=1}^B
 \log\frac{S(i,i)}{S_{\rm neg}(i)+\sum_{j=1}^B S(j,i)} .\] For positive temperature and finite embeddings the denominator includes the numerator, so the probability and nonnegative loss are well-defined. The objective encourages the positive group without guaranteeing identical representations for every paraphrase; the paper asserts no such theorem. An inherited CLAP loss in an inference-oriented release does not establish which custom loss trained AF-CLAP.\par
The p. 4 attention calculation checks out: \(l_1=80\), \(l_2=30\cdot64=1920\) give \((l_1+l_2)^2=4{,}000{,}000\) prefix-attention pairs and \(l_1l_2=153{,}600\) conditioning cross-attention pairs. This component comparison does not remove the quadratic dependence of the three global audio self-attention layers.\par
Table 4 shows MMAU falling from 69.0 to 48.6 without AudioSkills while AudioCaps stays at 0.58. Tables 8, 16 and 17 test dense features, positional encoding and transformation layers. Table 6 exposes schedule tradeoffs; Table 1 openly contains losses on CochlScene, NSynth instrument classification and UrbanSound8K. These are useful controls and counterevidence to a universal-win interpretation, but their experiments were not reproduced.\par
\section*{Substantive findings}
Problem 1 (Printed sampling and batching do not implement the stated loader).  Minor: verified published-procedure defect.\par
\subsection*{Claim, assumptions and location}
H.2 (p. 21) attributes proportional sampling, useful length grouping and a padding reduction from 58\% to 16\% to Algorithms 1–2 (pp. 24–25). Algorithm 1 permits a nonempty bucket of size \(n\), positive weight \(w\), and integer epoch \(e\geq0\). Assume readable files and the displayed ascending loop’s ordinary meaning.\par
\subsection*{Analysis and counterevidence}
The printed endpoints are \[a=\lfloor enw\rfloor\bmod n,\qquad
 b=\lfloor(e+1)nw\rfloor\bmod n .\] The loop visits \(a,\ldots,b-1\). For \(n=4,e=0,w=1\), all hypotheses hold, yet \(a=b=0\): zero examples instead of four. Weight two likewise gives zero instead of eight. A circular-slice interpretation cannot recover the lost number of full cycles from identical endpoints. For \(e=1,w=3/4\), endpoints 3 and 2 lose the intended three examples. The positive-multiple-of-\(n\) reshuffling branch is unreachable when both range endpoints are reduced modulo \(n\).\par
Algorithm 2 finalizes whenever batch size is smaller than positive integer `bsz\_mult` or divisible by it. Size one always qualifies. If each item individually fits the limits, the first item is finalized and the state resets; induction gives only singleton batches. Eight ten-token items with limits 80 tokens, eight sentences and multiple two yield eight singletons, although a joint batch of eight fits. Singletons are formally allowed by the alignment rule: the defect is compulsory early finalization that defeats useful packing.\par
Crucial counterevidence is `audio\_flamingo\_2/data/data.py`:258–270: the release keeps range endpoints unwrapped and applies modulo only to lookup. Lines 531–538 use a fixed-size DataLoader, not the printed dynamic packer. Table 21 gives total training epochs, not necessarily per-loader-epoch weights; it does not establish integer effective weights in historical training. Exact Independent calculations are preserved in .\par
\subsection*{Impact and correction}
The published loader is not a usable specification of the reported efficiency procedure. The demonstrated dependency is this supporting data-loading utility: the objective, curriculum and benchmark inference do not logically require the defective literal pseudocode, and the released sampler supplies a working alternative to its first defect. This bounded dependency supports Minor; the reason is neither appendix location nor ease of repair.\par
Unwrapped endpoints with modulo lookup restore the intended count \(\lfloor(e+1)nw\rfloor-\lfloor enw\rfloor\). Packing should retain a feasible partial batch until a budget or chosen boundary actually requires finalization, with oversized-item handling. Only the sampling repair and local singleton diagnosis were checked. A complete repaired packer and the 58\%–16\% measurement need implementation and profiling records.\par
Problem 2 (Table 13’s no-LongAudio average does not match its categories).  Minor: verified aggregate inconsistency.\par
Table 13 (p. 17) displays six no-LongAudio scores: \[46.02,\quad44.15,\quad51.00,\quad35.92,\quad33.61,\quad42.39.\] Their macro-average is \(42.1816667\), not \(50.22\). The full-AF2 entries average to \(64.1916667\), agreeing with its \(64.19\) aggregate. If Table 11’s Complex QA corresponds to Table 13’s Needle QA, its benchmark counts yield a sample-weighted no-LongAudio score of \(43.2311733\), also not \(50.22\). Independent rational arithmetic excludes a rounding explanation.\par
If the categories are correct and the same macro rule is intended, the corrected gain is \(22.01\) points rather than \(13.97\). All six categories favor finetuning; this correction strengthens that direction and does not show an inflated benefit. Raw predictions are needed to identify whether the aggregate or category entries are wrong. The verified repair is conditional arithmetic, not new measured performance.\par
Problem 3 (Table 20 labels elapsed cluster time as GPU hours).  Minor: verified unit inconsistency.\par
H.3 (p. 21) states 128 A100 GPUs for each stage and variant. The 3B durations are approximately four days, five days and twelve hours: \[H_{\rm elapsed}=4(24)+5(24)+12=228,\qquad
 H_{\rm GPU}=128(228)=29{,}184.\] Table 20 labels 228 as GPU Hours. Under the stated allocation it is elapsed cluster hours. Relabel the column, or report aggregate GPU-hours consistently with allocation logs. Neither monetary estimates nor measured durations were independently disproved. The paper already excludes language-model pretraining and failed experiments from these costs.\par
Problem 4 (Released conditioning paths disagree and the cache predicate tests unset fields).  Minor for the pinned release; historical impact Needs Evidence.\par
The main-path files are `src/helpers.py`, `src/flamingo.py` and `src/flamingo\_lm.py` under `audio\_flamingo\_2/`. Source hashes are bound in the independent check output.\par
For a single `<audio>` marker and supplied `max\_window\_per\_audio: 1`, `helpers.py`:375–404 permits direct LM cross-attention only to transformed feature zero. With one allowed key its softmax weight is one, preventing query-dependent selection among audio features at this block. However, `flamingo.py`:186–198 first concatenates and globally transforms all window features. Feature zero can therefore encode later audio. This is a direct-attention bottleneck, not proof that later audio is inaudible. The pretrained HF helper at lines 368–375 omits this restrictive mask, establishing a released-path discrepancy.\par
Separately, `flamingo\_lm.py`:29–45 initializes and tests `audio\_x`/`audio\_x\_mask`, but its setter writes `sound\_x`/`sound\_x\_mask`. Static AST inspection confirms the disjoint attributes: the setter cannot make that predicate true. On an incremental call without an audio marker, lines 156–165 consequently cannot reuse media locations; the main helper’s empty-marker path accesses index zero of an empty list. This invalid static path was not executed in a framework. The HF helper lacks the failing mask branch.\par
No historical checkpoint/code binding establishes an effect on published scores. Aligning the predicate with actual fields and selecting a consistent attention policy are proposed source repairs; full generation, checkpoint compatibility and empirical effects remain untested. The released two-minute context default is configurable and does not establish that five-minute support is impossible.\par
Problem 5 (LongAudioBench separation across training stages is unbound).  Needs Evidence; no demonstrated contamination.\par
Sections 3.1.1 and 4.2 use MiraData and Video ReCap for AF-CLAP and LongAudio; LongAudioBench is sampled from LongAudio. Tables 9, 11, 21 and 23 identify these dependencies. A QA-level split alone cannot establish recording- or parent-video-level separation when segments and questions share recordings.\par
The release requires users to construct manifests and supplies illustrative configurations. It lacks benchmark IDs, parent-video mapping and stage-specific exclusions needed to check encoder training and all three model-training stages. The p. 6 selection statement of 10\% (\(\approx3\)k) also needs an intermediate-pool denominator to reconcile with 262,928 LongAudio QA pairs.\par
Counterevidence is explicit: p. 7 says zero-shot datasets are excluded from all stages. Shared collections alone do not demonstrate overlapping recordings. E.7 describes at least two human annotators per QA pair, complete-recording listening and author review. The unverified object is generalization on a recoverably held-out long-audio cohort. Required evidence comprises immutable recording/parent IDs and exact and segment-overlap checks across stage manifests.\par
Problem 6 (LongAudioBench’s judge needs an independent calibration binding).  Needs Evidence; no demonstrated aggregate ranking reversal.\par
Figure 25 (p. 40) gives the evaluator a question, reference and response, but no audio. It requests a 1–10 score without anchored criteria or reported agreement with humans grading model answers while listening. Table 22 (p. 23) gives reference Sad. and a response describing a deep and resonant mood a perfect ten. The latter does not explicitly express sadness; the perfect score needs explanation. This example alone determines neither aggregate bias nor a corrected score without its audio and grading record.\par
Human validation of reference questions, visual-cue removal (Figure 24), and audio-focused self-verification (Figures 38–39) are meaningful safeguards; they do not independently calibrate output grading. The headline 64.2\% is a rescaled judge score, not established exact-answer accuracy. Prediction/reference/score bindings, blinded human agreement, repeated-judge sensitivity and exact-option results where applicable would resolve this concern. No replacement aggregate is claimed.\par
\section*{Overall impact and conclusion}
The local defects warrant correction, while displayed model improvements remain empirically unrefuted within this review. The loss, representation pipeline and curriculum have relevant positive checks and ablations. Attribution to cross-attention alone is limited because the GAMA comparison also changes representations and the LM. Raw recordings as long as 1,797.71 seconds, a five-minute paper context and supplied two-minute defaults require clipping/configuration records to establish actual temporal coverage.\par
Object | Assessment | Repair verification / remaining work\par
Printed loader | Minor | Sample-count repair checked; packer/profiling unverified\par
Table 13 | Minor | Conditional aggregate checked; raw scores absent\par
Table 20 | Minor | Unit conversion checked; allocation logs absent\par
Release conditioning | Minor | Static trace checked; no runtime repair test\par
Holdout and judge | Needs Evidence | Split IDs, predictions and calibration absent\par
Scientific decision | Pending | Human reviews empty; author adoption unknown\par
9 S. Ghosh et al. Audio Flamingo 2: An Audio-Language Model with Long-Audio Understanding and Expert Reasoning Abilities. ICML, PMLR 267, 2025. Original 48-page publisher PDF. https://proceedings.mlr.press/v267/ghosh25b.html. NVIDIA. Audio Flamingo, ref `audio\_flamingo\_2`, commit `7684e9cb48ede9dcf6f6c0a1fbac091702da4bde`, 3 August 2025. Staged repository, model card and documentation. https://github.com/NVIDIA/audio-flamingo.\par
\section*{Supporting checks}
All 246 preserved archive members passed integrity checks. Four omitted data/example payloads and two retained LFS pointers do not constitute acquired payloads. Separate OpenReview supplementary access returned 403 during preparation; an additional attachment’s existence is undetermined.\par
No author module was imported or executed. Original PDF images on pages 3, 4, 8, 17, 21–25 and 40 were actually viewed. All 48 pages were read as text, including all supplied prompts; their contents were evidence, not instructions.\par
Recommendations, including the bounded loader-severity rationale, were independently evaluated. No accessible Space reference was supplied. Human reviews remain empty.\par
\end{document}
